\subsection{RANK OF A LINEAR MAPPING}

DEFINITION 3. Let E, F be two vector spaces over a field K. For every linear mapping u of E into F, the dimension of the subspace u(E) of F is called the rank of u and denoted by rg(u).

If \(\mathbf{N} = \operatorname{Ker}(u)\) , \(\mathbf{E} / \mathbf{N}\) is isomorphic to \(u(\mathbf{E})\) , whence the relation

\[
\operatorname{rg} (u) = \operatorname{codim} _ {\mathrm{E}} (\operatorname{Ker} (u))\tag{11}
\]

and therefore

\[
\operatorname{rg} (u) + \dim (\operatorname{Ker} (u) = \dim E.\tag{12}
\]

Moreover, by formula (3)

\[
\operatorname{rg} (u) + \dim (\operatorname{Coker} (u)) = \dim F.\tag{13}
\]

PROPOSITION 9. Let E, F be two vector spaces over a field K and u:E→F a linear mapping.

\begin{enumerate}
\def\labelenumi{(\roman{enumi})}
\item
  \(\operatorname{rg}(u) \leqslant \inf(\dim E, \dim F)\) .
\item
  Suppose that \(\mathbf{E}\) is finite-dimensional; in order that \(\operatorname{rg}(u) = \dim \mathbf{E}\) , it is necessary and sufficient that \(u\) be injective.
\item
  Suppose that \(\mathbf{F}\) is finite-dimensional; in order that \(\operatorname{rg}(u) = \dim \mathbf{F}\) , it is necessary and sufficient that \(u\) be surjective.
\end{enumerate}

This follows immediately from relations (12) and (13).

COROLLARY. Let E be a vector space of finite dimension n and u an endomorphism of E. The following properties are equivalent:

\begin{enumerate}
\def\labelenumi{(\alph{enumi})}
\item
  \(u\) is bijective;
\item
  \(u\) is injective;
\item
  \(u\) is surjective;
\item
  \(u\) is right invertible;
\item
  \(u\) is left invertible;
\item
  \(u\) is of rank \(n\) .
\end{enumerate}

If E is an infinite-dimensional vector space, there are injective (resp. surjective) endomorphisms of E which are not bijective (Exercise 9).

Let K, \(K'\) be two fields, \(\sigma:K\to K'\) an isomorphism of K onto \(K'\) , E a vector K-space, \(E'\) a vector \(K'\) -space and \(u:K\to K'\) a semi-linear mapping relative to \(\sigma\) (§1, no. 13); the dimension of the subspace \(u(\mathbf{E})\) of \(E'\) is also called the rank of u. It is also the rank of u considered as a linear mapping of E into \(\sigma_{*}(\mathbf{E}')\) , for every basis of \(u(\mathbf{E})\) is also a basis of \(\sigma_{*}(u(\mathbf{E}))\) .

